\documentclass[10pt, letterpaper]{article}

% Packages
\usepackage[
ignoreheadfoot,
top=1.5 cm,
bottom=1.5 cm,
left=2 cm,
right=2 cm,
footskip=1.0 cm
]{geometry}

\usepackage{titlesec}
\usepackage{tabularx}
\usepackage{array}
\usepackage[dvipsnames]{xcolor}
\definecolor{primaryColor}{RGB}{0,0,0}
\usepackage{enumitem}
\usepackage{fontawesome5}
\usepackage{amsmath}

\usepackage[
pdftitle={Luis Antonio Flores Esquivel CV},
pdfauthor={Luis Antonio Flores Esquivel},
colorlinks=true,
urlcolor=primaryColor
]{hyperref}

\usepackage{bookmark}
\usepackage{changepage}
\usepackage{paracol}
\usepackage{needspace}
\usepackage{iftex}


% ATS-readable PDF
\ifPDFTeX
\input{glyphtounicode}
\pdfgentounicode=1
\usepackage[T1]{fontenc}
\usepackage[utf8]{inputenc}
\usepackage{lmodern}
\fi

\usepackage{charter}

% Layout settings
\raggedright
\pagestyle{empty}
\setcounter{secnumdepth}{0}
\setlength{\parindent}{0pt}
\setlength{\columnsep}{0.15cm}

\titleformat{\section}{\needspace{4\baselineskip}\bfseries\large}{}{0pt}{}[\titlerule]
\titlespacing*{\section}{0pt}{0.7\baselineskip}{0.5\baselineskip}

\renewcommand\labelitemi{$\vcenter{\hbox{\small$\bullet$}}$}

% Custom environments
\newenvironment{highlights}{
	\begin{itemize}[topsep=0.15cm,itemsep=0pt,leftmargin=10pt]
	}{
	\end{itemize}
}

\newenvironment{onecolentry}{
	\begin{adjustwidth}{0cm}{0cm}
	}{
	\end{adjustwidth}
	\vspace{0.02cm}
}

\newenvironment{twocolentry}[1]{
	\onecolentry
	\columnratio{0.75} 
	\begin{paracol}{2}
		\raggedright #1
		\switchcolumn
	}{
	\end{paracol}
	\endonecolentry
}


\newenvironment{header}{
	\centering\linespread{1.4}
}{}

% ================= DOCUMENT =================
\begin{document}
	
	\begin{header}
		{\fontsize{20pt}{20pt}\selectfont Luis Antonio Flores Esquivel}
		
		\vspace{6pt}
		\makebox[\textwidth][c]{%
			\faEnvelope\,
			\href{mailto:flores.esquivel.luis.antonio2@gmail.com}{flores.esquivel.luis.antonio2@gmail.com}
			\;|\;
			\faPhone*\,
			+52 55 5102 2244
			\;|\;
			\faLinkedin\,
			LinkedIn: \href{https://linkedin.com/in/lflorese}{lflorese}
			\;|\;
			\faGithub\,
			GitHub: \href{https://github.com/LuisFlorEsq}{LuisFlorEsq}
		}

	\end{header}

	
	% ================= EDUCACIÓN =================
	\section{Educación}
	
	\begin{twocolentry}{\textbf{Escuela Superior de Cómputo (ESCOM), Instituto Politécnico Nacional (IPN)}\\
	Licenciatura en Ingeniería en Inteligencia Artificial}
	{Graduación: Enero 2026}

	\end{twocolentry}
	
	\begin{onecolentry}
		\begin{highlights}
			\item Cursos relevantes: Aprendizaje Automático, Aprendizaje Profundo, Procesamiento de Lenguaje Natural (PLN), Visión por Computadora, Computación Bioinspirada
		\end{highlights}
	\end{onecolentry}
	
	% ================= EXPERIENCIA =================
	\section{Experiencia Profesional}
	
	\begin{twocolentry}
		{\textbf{IT Intern - Inteligencia Artificial}
		\\ El Puerto de Liverpool S.A.B. de C.V., Santa Fe, Ciudad de México}
		{Abr 2025 -- Sept 2025}
	\end{twocolentry}
	
	\begin{onecolentry}
		\begin{highlights}
			\item Diseño, desarrollo y despliegue de soluciones de inteligencia artificial en producción para la automatización de procesos en áreas de comercio electrónico, logística y legal.
			
			\item Desarrollo de microservicios escalables utilizando FastAPI, Docker y Google Cloud Platform (GCP).
			
			\item Implementación de modelos de visión por computadora (YOLO, DETR, U-Net) para detección de objetos y edición automática de imágenes para el foro fotográfico de Liverpool.
			
			\item Integración de modelos generativos (Google Gemini, Imagen) para generación y edición automatizada.
			
			\item Implementación de algoritmos de optimización para ruteo logístico y estimación de tiempos de entrega.
			
			\item Colaboración en la propuesta y desarrollo de un sistema RAG sobre documentos del área legal.
			
			\item Optimización de consultas SQL en BigQuery mediante estrategias de caching y mejora del rendimiento.
			
		\end{highlights}
	\end{onecolentry}
	
	% ================= PROYECTOS =================
	\section{Proyectos relevantes}
	
	\begin{twocolentry}
		{\textbf{Sistema de Autenticación Biométrica mediante Reconocimiento Facial\\
			(Proyecto de Titulación)}\\
		FastAPI, Docker, FaceNet, Microsoft Azure, Pinecone}
		{Ago 2024 -- Jun 2025}
	\end{twocolentry}
	
	\begin{onecolentry}
		\begin{highlights}
			\item Desarrollo de un sistema de autenticación biométrica basado en reconocimiento facial y Deep Learning.
			\item Desarrollo de APIs REST con FastAPI (Python) para consumo desde una aplicación móvil en Kotlin.
			\item Uso de bases de datos vectoriales (Pinecone) para almacenamiento de embeddings faciales y bases de datos relacionales para gestión de usuarios.
			\item Despliegue de servicios utilizando Docker y Microsoft Azure.
		\end{highlights}
	\end{onecolentry}
	
	\vspace{0.2cm}
	
	\begin{twocolentry}
		{\textbf{DOF – Inteligencia Artificial (Diario Oficial de la Federación)}\\
		Django, LangChain, Hugging Face, OpenAI API, Pinecone}
		{Mar 2024 -- Jun 2024}
	\end{twocolentry}
	
	\begin{onecolentry}
		\begin{highlights}
			\item Desarrollo de una aplicación web basada en LLMs para consulta y resumen de documentos gubernamentales.
			\item Evaluación de modelos de lenguaje de código abierto y propietarios para tareas de Question Answering y Retrieval-Augmented Generation (RAG).
			\item Implementación de un chatbot conversacional con LangChain y Pinecone.
			\item Desarrollo de servicios backend con Django para interacción en tiempo real.
		\end{highlights}
	\end{onecolentry}

	
	% ================= HABILIDADES =================
	\section{Habilidades Técnicas}
	
	\begin{onecolentry}
		\textbf{Lenguajes de Programación:} Python, Java, C, C++, MATLAB
	\end{onecolentry}
	
	\begin{onecolentry}
		\textbf{Frameworks y Backend:} FastAPI, Django, Flask, LangChain
	\end{onecolentry}
	
	\begin{onecolentry}
		\textbf{Machine Learning y Deep Learning:} Scikit-learn, TensorFlow, PyTorch, NumPy, Pandas, OpenCV
	\end{onecolentry}
	
	\begin{onecolentry}
		\textbf{Cloud y DevOps:} Docker, Google Cloud Platform (GCP), Microsoft Azure, Git, GitHub
	\end{onecolentry}
	
	\begin{onecolentry}
		\textbf{Bases de Datos:} Bases de datos vectoriales (Pinecone, Weavite) Bases de datos relacionales (MySQL, PostgreSQL)
	\end{onecolentry}
	
	% ================= CERTIFICACIONES =================
	\section{Certificaciones y Formación Complementaria}
	
	\begin{onecolentry}
		\begin{highlights}
			\item Google Cloud Computing Foundations Certificate – Google Cloud Platform
			\item Escuela de Verano de Ciencias Cognitivas Computacionales y Procesamiento de Lenguaje Natural (EVCCCPLN 2023, 2024)
		\end{highlights}
	\end{onecolentry}
	
\end{document}
